\documentclass[12pt]{article}
\usepackage[margin=0.65in]{geometry}

\newcommand{\duedate}{02/24/2026}
\newcommand{\assignment}{Continual Multimodal Pretraining}

\newcommand{\name}{Aksel Kretsinger-Walters, adk2164}
\newcommand{\email}{adk2164@columbia.edu}

\makeatletter
\def\input@path{{../}{../../}{../../../}}
\makeatother
\input{pset_template_CLMM.tex} %% DO NOT CHANGE THIS LINE

\begin{document}
\psetheader %% DO NOT CHANGE THIS LINE

\section*{A Practitioner's Guide to Continual Multimodal Pretraining}

\paragraph{Motivation.}
Multimodal foundation models (e.g., vision-language models) are typically pretrained on massive static
datasets. In practice, however, data evolves over time and arrives sequentially. Continual multimodal
pretraining (CMP) studies how to update multimodal models without degrading prior capabilities.

\paragraph{Core Challenges.}
Continual multimodal learning faces unique issues:
\begin{itemize}
		\item Cross-modal alignment drift (vision and language embeddings desynchronize).
		\item Catastrophic forgetting of earlier modalities or domains.
		\item Replay inefficiency at multimodal scale.
\end{itemize}

\paragraph{Empirical Study.}
The paper systematically evaluates continual pretraining strategies on vision-language models under domain
shifts and modality imbalance.

Methods analyzed include:
\begin{itemize}
		\item Finetuning
		\item Rehearsal-based replay
		\item Regularization methods
		\item Parameter-efficient tuning
\end{itemize}

\paragraph{Key Observations.}
\begin{itemize}
		\item Naive finetuning severely degrades zero-shot transfer.
		\item Replay is critical but must preserve cross-modal alignment.
		\item Parameter-efficient methods mitigate forgetting while controlling compute.
		\item Continual multimodal pretraining is substantially more fragile than unimodal CL.
\end{itemize}

\paragraph{Practical Recommendations.}
\begin{itemize}
		\item Maintain balanced replay across modalities.
		\item Monitor alignment metrics, not just downstream accuracy.
		\item Use lightweight adaptation modules rather than full-model updates.
\end{itemize}

\paragraph{Takeaway.}
Continual multimodal pretraining is not merely “CL + multimodality.” Cross-modal coupling introduces new
failure modes. Stable long-term multimodal learning requires explicit mechanisms to preserve representational
alignment across modalities and time.

\end{document}
